\noindent
{\Large \textbf{\section*{Заключение}}} \par
\addcontentsline{toc}{section}{Заключение} \par

В ходе выполнения лабораторной работы реализован двумерный клеточный автомат с окрестностью фон Неймана и тороидальными границами по варианту $N=50045688$. 

Реализовано:
\begin{enumerate}[itemsep=-3pt]
 \item Правило перехода в соответствии с номером варианта (таблица истинности булевой функции 5 переменных)
 \item Поле с тороидальными граничными условиями
 \item Возможность изменения начальных условий: ширины поля, числа итераций, типа генерации, режима работы (автоматический или пошаговый)
\end{enumerate}

\textbf{Достоинства реализованного решения:}
\begin{enumerate}[itemsep=-3pt]
 \item Правило перехода строится из исходного числа автоматически, поэтому для другого варианта достаточно изменить одну константу.
 \item Пользователю предоставлено два режима работы программы: с заданным числом итераций и пошаговый.
 \item Таблица истинности строится один раз в конструкторе и хранится в ассоциативном массиве из 32 элементов, поэтому при вычислении шага для клетки нужен лишь поиск по ключу (не более $\log_2 32 = 5$ сравнений), а не вычисление значения булевой функции.
\end{enumerate}

\textbf{Недостатки и возможности для улучшения:}
\begin{enumerate}[itemsep=-3pt]
 \item Номер текущей итерации хранится в программе, но не отображается в интерфейсе, также нет вывода статистики (например, доли живых клеток). Для этого можно реализовать отдельный виджет со статистикой.
 \item Кнопка «Перегенерировать» не останавливает таймер. Если нажать её во время автоматического режима, оставшиеся итерации продолжат выполняться уже на новом поле. Решение - останавливать таймер и сбрасывать счетчик оставшихся шагов в \texttt{on\_init}
\end{enumerate}

\textbf{Масштабируемость решения:}
\begin{enumerate}[itemsep=-3pt]
    \item Добавить возможность сохранять и загружать файлы с состоянием поля.
    \item Добавить поддержку правил, основанных не только на окрестности фон Неймана, но и на окрестности Мура.
    \item Добавить возможность делать шаг не только вперед ($t+1$), но и назад ($t-1$).
\end{enumerate}

\newpage
